\section{The Act, the Code, and the Book}

Three documents make the calendar of the 1962 \latin{Missale Romanum} what it is, and they
are not interchangeable. An apostolic letter given \latin{motu proprio} approved a new body of
rubrics; a general decree promulgated it; and a typical edition of the Missal printed it and
put it to work. Every rule in this reference is stated at its own place in one of them, and
where those places disagree — as at four points they measurably do — the disagreement is
reported rather than smoothed.

\subsection{What the 1960 reform was for}

John XXIII's \latin{Rubricarum instructum} of 25 July 1960 opens with a diagnosis rather than
a programme. The Apostolic See, it says, has continuously worked since Trent to define and
order the body of rubrics governing public worship more exactly; the result was accretion:

\begin{studyframe}
\latin{Pluribus itaque emendationibus, variationibus et additamentis decursu temporis
introductis, totum rubricarum systema abunde succrevit, non semper vero systematico ordine
servato, et non sine primitivae perspicuitatis ac simplicitatis detrimento.}
\medskip

\textit{Working gloss.} With many emendations, variations, and additions introduced in the
course of time, the whole system of rubrics grew abundantly, though not always with a
systematic order preserved, and not without loss of its original clarity and simplicity.
\end{studyframe}

The letter then narrates the immediate prehistory in two steps. Pius XII, acceding to the
petitions of many bishops, had judged that the rubrics of the Breviary and Missal should be
reduced in certain respects to a simpler form, and this was done by the general decree of the
Sacred Congregation of Rites of 23 March 1955. In the following year, with preparatory studies
for a general liturgical restoration maturing, the same pope consulted the bishops on the
liturgical emendation of the Breviary and, having weighed their responses, judged that a
general and systematic emendation of the rubrics of both books should be undertaken, entrusting
it to the commission of experts already charged with the studies for the general restoration.

The decisive paragraph is John XXIII's own. Having determined to convoke an ecumenical council,
he considered more than once what should be done about his predecessor's undertaking, and came
to this view: \latin{altiora principia, generalem liturgicam instaurationem respicientia, in
proximo Concilio Oecumenico Patribus esse proponenda; memoratam vero rubricarum Breviarii ac
Missalis emendationem diutius non esse protrahendam} — the higher principles bearing on a
general liturgical restoration were to be proposed to the Fathers in the coming Council, but the
emendation of the rubrics of the Breviary and Missal was not to be delayed any longer. The 1960
code is therefore, by its own account, a deliberately interim act: a tidying of the existing law
undertaken precisely because the substantial question had been reserved to a council that had
not yet met. A reader who treats the 1962 calendar as the settled endpoint of a tradition is
reading against the promulgating document; so is a reader who treats it as the first move of
the reform that followed.

\subsection{The six dispositions}

The operative part of \latin{Rubricarum instructum} is six numbered paragraphs, and their
content controls a good deal of what follows in this reference.

\begin{studybox}{\latin{Rubricarum instructum}, nn. 1--6, in substance}
\begin{enumerate}[leftmargin=1.5em, itemsep=0.25em]
\item The new code of rubrics of the Roman Breviary and Missal, \latin{per tres partes
digestum} — arranged in three parts, namely \latin{Rubricae generales}, \latin{Rubricae
generales Breviarii Romani}, and \latin{Rubricae generales Missalis Romani} — to be published
shortly by the Sacred Congregation of Rites, is to be observed from 1 January 1961 by all who
follow the Roman Rite. Those who observe another Latin rite must conform themselves as soon as
possible to the new code and to the Calendar in everything not strictly proper to that rite.
\item On the same 1 January 1961 the former \latin{Rubricae generales Breviarii et Missalis
Romani} and the \latin{Additiones et Variationes} made under Pius X's bull \latin{Divino
afflatu} cease to have force, as does the general decree of 23 March 1955 on reducing the
rubrics to a simpler form, which has been taken up into the new redaction. Decrees and replies
to doubts of the same Congregation that do not agree with the new form of the rubrics are
abrogated.
\item Statutes, privileges, indults, and customs of every kind that stand against these rubrics
are revoked — \latin{etiam saecularia et immemorabilia, immo specialissima atque individua
mentione digna}, even centuries-old and immemorial ones, indeed those deserving most special
and individual mention.
\item Approved publishers of liturgical books may prepare new editions of the Breviary and
Missal arranged according to the new code; the Congregation is to issue particular instructions
to secure their uniformity.
\item In new editions the old rubrical texts are omitted and the new ones prefixed: to the
Breviary, the \latin{Rubricae generales} and the \latin{Rubricae generales Breviarii Romani};
to the Missal, the \latin{Rubricae generales} and the \latin{Rubricae generales Missalis
Romani}.
\item All whom it concerns are to see that diocesan and religious calendars and propers are
conformed to the norm and mind of the new redaction of the rubrics and of the Calendar, to be
approved by the Sacred Congregation of Rites.
\end{enumerate}
\end{studybox}

Three consequences deserve to be stated plainly, because each is a recurring source of
confusion. First, the code binds from 1 January 1961, not from 1962: the 1962 Missal is a book
made to fit a law already in force, not the instrument that enacted it. Second, \rn{5} is why
the 1962 Missal prints parts one and three of the code and not part two — the Breviary rubrics
belong to the Breviary. A reader looking in the Missal for the rule that governs the Office
alone will not find it there, and this reference marks those places. Third, \rn{3} is the
reason a 1962 calendar cannot be reconstructed from older privileges: an immemorial local
octave, a proper rank, a customary commemoration, if it stands against the new rubrics, was
revoked, and nothing but a fresh concession from the Holy See restores it.

\subsection{The promulgating decree and what it added}

The Sacred Congregation of Rites promulgated the code by general decree on 26 July 1960, the
day after the papal approval. The decree is short and does one thing beyond promulgation:

\begin{studyframe}
\latin{Ut autem libri liturgici qui in usu habentur adhuc adhiberi possint, rubricarum codici
adduntur ``Variationes'' Breviariis et Missalibus necnon Martyrologio aptandis.}
\medskip

\textit{Working gloss.} But so that the liturgical books in use may still be usable, there are
added to the code of rubrics ``Variations'' for adapting Breviaries and Missals and also the
Martyrology.
\end{studyframe}

Those \latin{Variationes} matter far beyond their stated purpose. Their first chapter,
\latin{Variationes in Calendario}, is the only place where the 1960 reform states its own
calendar changes as changes — which feasts were regraded, which were reduced to
commemorations, which were struck from the calendar, which were newly inscribed, which were
moved, and which were renamed. The 1962 Missal simply prints the result. Anyone who wants to
know what 1960 did, as distinct from what 1962 shows, must read the \latin{Variationes};
Section~\ref{sec:sanctoral} of this reference does so.

The same decree explains a fact about the Acta text that a reader will otherwise find puzzling:
the code was printed in \latin{Acta Apostolicae Sedis} 52 (1960) at pages 597--740 together
with a complete new \latin{Calendarium Breviarii et Missalis Romani} at pages 686 and following,
and with the tables of occurrence and concurrence and their \latin{Notanda} to page 705, before
the \latin{Variationes} begin at page 706. The Acta therefore carry both the law and the
calendar it presupposes.

\subsection{The 1962 typical edition}

The Missal's own opening decree, dated at the offices of the Sacred Congregation of Rites on
23 June 1962 and subscribed by Cardinal Larraona as Prefect and Enrico Dante as Secretary,
states its logic in a single sentence: a new body of rubrics having been approved and
promulgated, \latin{vix fieri non potuit quin eadem Sacra Rituum Congregatio novam Missalis
romani editionem, ad dictum rubricarum codicem plane accommodatam, pararet} — it could scarcely
not happen that the same Congregation should prepare a new edition of the Roman Missal fully
accommodated to that code of rubrics. The Vatican press printed it; the Congregation reviewed
it; and the decree declares this Vatican edition \latin{typica}, to be reproduced by those
competent \latin{in omnibus et per omnia, integre et absolute}.

That decree contains one of the four textual divergences this reference reports. It dates
\latin{Rubricarum instructum} to \latin{die 23 iulii anno 1960}. The motu proprio itself,
in the Acta, is dated \latin{die 25 mensis Iulii, anno 1960}; and the same Missal decree
continues \latin{posteroque die a Sacra Rituum Congregatione promulgato}, on the following day
promulgated by the Sacred Congregation of Rites, while the promulgating decree is dated 26 July.
The Missal's own two statements are therefore inconsistent with each other: 23 July plus one day
is not 26 July. The Acta dates are the ones this reference uses, and the Missal's ``23 iulii''
is treated as a printing error. Section~\ref{sec:collation} sets out the evidence.

\subsection{How to read a rule in this reference}

The code numbers its paragraphs continuously through all three parts: \rnn{1--137} are the
\latin{Rubricae generales}, \rnn{138--268} the Breviary rubrics, and \rnn{269--530} the Missal
rubrics. Citations here give the number alone for part one, and say ``Missal rubrics'' when the
number belongs to part three. Where the Missal prints something the code does not — a rubric at
the head or foot of a formulary, a note in the front matter, an entry in the calendar — the
citation names the Missal.

One structural rule of the code governs the whole of what follows and should be fixed before
anything else. \rn{7} reads: \latin{Praecedentia inter singulos dies liturgicos unice per
peculiarem tabellam (n. 91) determinatur} — precedence among individual liturgical days is
determined \emph{solely} by the particular table at \rn{91}. \rn{91} repeats the exclusivity
from the other side: precedence is governed by the following table \latin{sublatis quibuslibet
aliis titulis vel normis}, all other titles or norms being removed. The 1960 system is
therefore not a system of dignities from which precedence may be inferred. It is a system in
which one twenty-eight-line table decides, and in which the classes are labels for positions in
that table rather than the reverse.
